\section{The Spiritual Writer and the Last Years (1607--1621)}

\subsection{From controversy to the Psalms}

The last fifteen years in Rome were outwardly a rhythm of
congregations---Holy Office, Index, Rites, and the rest---and
inwardly a turn of the pen. In 1611 appeared the commentary
\work{In omnes Psalmos dilucida expositio}, of which the 1931
doctoral letter says simply that in explaining the Psalter ``he
joined learning with piety'' (AAS~23, 436; M). From 1615 to 1620,
in the octave-long retreat he made each year, he wrote the five
small books on which his devotional fame rests; the 1907 Catholic
Encyclopedia lists them with their dates, noting they were
``spiritual works written during his annual retreats.'' Alongside
them stand the memoir of 1613 and the late works of counsel and
controversy. The table gathers the writings named in the witnesses
used for this study; it is a reader's orientation, not a
bibliography of the whole corpus, which fills volumes of the old
\work{Opera omnia}.

\begin{fourstable}{Date}{Work}{Kind}{Witness for date}
\properrow{Louvain years}{\work{Institutiones linguae hebraicae}}
{Hebrew grammar from the Louvain teaching}{AAS 23, 434 (composed
young); print date not verified here}
\properrow{1586, 1588, 1593}{\work{Disputationes de controversiis
christianae fidei} (Ingolstadt, 3 tomes; later editions in 4)}{The
controversies}{AAS 23, 434; CE 1907}
\properrow{1597; 1598}{\work{Dottrina cristiana breve};
\work{Dichiarazione pi\`u copiosa della dottrina cristiana}}
{Catechisms, by order of Clement VIII}{CE 1907; Benedict XVI 2011;
AAS 23, 436}
\properrow{1604}{\work{Dichiarazione del Simbolo}}{Catechetical
exposition of the Creed}{CE 1907}
\properrow{1608}{\work{Responsio Matthaei Torti}}{Against the Oath
of Allegiance (pseudonymous)}{CE 1907}
\properrow{1609}{\work{Apologia pro responsione sua}}{Against
James I's \work{Premonition}, over his own name}{CE 1907}
\properrow{1610}{\work{Tractatus de potestate Summi Pontificis in
rebus temporalibus}}{Against Barclay; the indirect power}{CE 1907}
\properrow{1611}{\work{In omnes Psalmos dilucida expositio}}
{Commentary on the Psalms}{CE 1907; AAS 23, 436}
\properrow{1612}{\work{Admonitio ad Episcopum Theanensem}}{Counsel
to his nephew the bishop of Teano}{CE 1907; AAS 23, 436}
\properrow{1612}{\work{De scriptoribus ecclesiasticis}}{Patristic
reference chronology, from his lecture preparations}
{D\"ollinger--Reusch commentary (first Rome printing 1612); AAS 23,
434}
\properrow{1613}{\work{Vita} (autobiography)}{Third-person memoir,
written in June 1613 at Jesuit request}{the memoir's closing
sentence (A); CE 1907}
\properrow{1615}{\work{De ascensione mentis in Deum per scalas
rerum creatarum} (The Mind's Ascent to God by the Ladder of Created
Things)}{Retreat book}{CE 1907; Benedict XVI 2011}
\properrow{1615}{\work{Conciones habitae Lovanii}}{The Louvain
sermons, published late}{CE 1907}
\properrow{1616}{\work{De aeterna felicitate sanctorum}}{Retreat
book}{CE 1907}
\properrow{1617}{\work{De gemitu columbae} (The Mourning of the
Dove)}{Retreat book: the Church's penitential reform of itself}
{CE 1907; Benedict XVI 2011}
\properrow{1618}{\work{De septem verbis Christi}}{Retreat book: the
Seven Last Words}{CE 1907}
\properrow{1620}{\work{De arte bene moriendi} (The Art of Dying
Well)}{Retreat book}{CE 1907; Benedict XVI 2011}
\end{fourstable}

\noindent The autobiography's own late-works inventory---Psalms
commentary, the Italian tracts against the Venetian doctors, the
book against the King of England, against Barclay, against
Widdrington, and \work{De scriptoribus ecclesiasticis} ``with a
chronology''---confirms the list's core from the subject's side
(D\"ollinger--Reusch, 43; A).

\subsection{The Mind's Ascent and the Art of Dying}

The retreat books made Bellarmine a master of the interior life
for readers who never opened the \work{Disputationes}. The
\work{De ascensione mentis in Deum}, built consciously on the
Bonaventuran itinerary of the mind's ascent through creatures,
was the most loved; Benedict~XVI's 2011 catechesis quotes its
kernel: ``Whoever finds God finds everything, whoever loses God
loses everything'' (M, reported, in the Vatican's English). The
\work{De arte bene moriendi} distilled the tradition's rules into
two plain propositions---live well, and rehearse death while
living, by regular meditation and by laying up treasure in heaven
(M, reported). The \work{De gemitu columbae} turned the
controversialist's weapons around: it is an appeal for the reform
not of the heretics but of the Church's own clergy and people
(M, reported). These books circulated in dozens of editions and
translations across confessional lines---an irony their author
lived to see, and one the reception section takes up.

\subsection{Death at Sant'Andrea}

His health failing, Bellarmine obtained leave from Gregory~XV in
the late summer of 1621 to retire to the Jesuit novitiate of
Sant'Andrea on the Quirinal, ``to prepare for death, which he felt
near'' (\work{Lux illa}, AAS~22, 596; M; CE 1907 concurs). The
decretal's account of the deathbed is the account of the
canonization process: patient endurance, the last sacraments, the
visits of pope and cardinals, and the crowds. He died on
17~September 1621, aged seventy-eight---on the feast, the decretal
notes, of the Stigmata of St.~Francis, a commemoration whose
universal celebration he himself had promoted (AAS~22, 596; M).
The funeral overwhelmed the plans: though his will had asked for
burial without pomp, Gregory~XV ordered exequies at his own
expense; the crowd pressed to touch the body, to kiss hands and
feet, to carry off threads of the mitre and of the cardinal's hat
as relics; soldiers around the bier could hardly hold the press,
and the people ``with one voice acclaimed Robert a saint''
(AAS~22, 596; M, resting on process testimony; E in character).
He was buried in the Ges\`u, first in the common sepulcher of the
Society's priests, then in the crypt where Ignatius himself had
lain (AAS~22, 596; M). His own dearest wish about burial---to lie
near Aloysius Gonzaga---waited three hundred and two years.
